% concatenated from counterexample_RC.tex
\documentclass[11pt]{article}

\usepackage[margin=1in]{geometry}
\usepackage{amsmath,amssymb,amsthm}
\usepackage{booktabs,array}
\usepackage{enumitem}
\usepackage{hyperref}
\usepackage{microtype}

\newtheorem{theorem}{Theorem}
\newtheorem{proposition}[theorem]{Proposition}
\newtheorem{lemma}[theorem]{Lemma}
\newtheorem{remark}[theorem]{Remark}

\theoremstyle{definition}
\newtheorem{definition}{Definition}[section]

\newcommand{\forb}{\operatorname{forb}}
\newcommand{\Avoid}{\operatorname{Avoid}}
\newcommand{\ex}{\operatorname{ex}}
\newcommand{\cF}{\mathcal F}

\title{A Counterexample to the Anstee--Sali Conjecture}
\author{Pei Wu}
\date{7 August 2026}

\begin{document}
	\maketitle
	
	\begin{abstract}
		This note presents a counterexample to the Anstee--Sali conjecture for forbidden
		configurations. The candidate is the $4$-uniform family
		\[
		\cF_2=\{xyab,xybc,xycd,xyda\}
		\]
		on six vertices: a fixed two-vertex core $\{x,y\}$ joined to the four edges of a
		$4$-cycle. We give an explicit certificate that every four-fold product whose
		factors are of type $I$, $I^c$, or $T$ contains $\cF_2$, whereas $I^3$ avoids it.
		Thus $X(\cF_2)=4$, so the conjecture predicts
		\[
		\forb(m,\cF_2)=\Theta(m^3).
		\]
		On the other hand, a result of Mubayi on complete multipartite hypergraphs
		implies
		\[
		\forb(m,\cF_2)=\Omega(m^{7/2}),
		\]
		which is asymptotically larger than $m^3$.
	\end{abstract}
	
	\section{Background and notation}
	
	In extremal graph theory, Tur\'an's theorem~\cite{Turan1941} determines the
	maximum number of edges in an $n$-vertex graph containing no copy of $K_r$,
	$r>2$: the extremal number is
	\[
	\left(\frac{r-2}{r-1}+o(1)\right)\binom{n}{2}.
	\]
	More generally, for a fixed graph $H$, the Tur\'an number $\ex(n,H)$ is the
	maximum number of edges in an $n$-vertex graph containing no copy of $H$.
	Forbidden-configuration problems ask an analogous question for set systems.
	
	The Sauer--Shelah lemma, equivalently the classical VC-dimension
	bound~\cite{Sauer1972,Shelah1972,VapnikChervonenkis1971}, implies that
	$\forb(m,F)$ is polynomially bounded in $m$ for every fixed nontrivial
	configuration $F$. A central conjecture of Anstee and Sali predicts that the
	correct polynomial order is determined by products of three standard set
	systems; see~\cite{AnsteeSali2013} for a survey.
	
	\subsection{Forbidden configurations of set systems}
	
	Let $X$ be a finite ground set. A \emph{set system} on $X$ is a family
	$\mathcal S\subseteq \mathcal P(X)$. If $Y\subseteq X$, the restriction of
	$\mathcal S$ to $Y$ is
	\[
	\mathcal S|_Y=\{A\cap Y:A\in\mathcal S\}.
	\]
	Two set systems are \emph{isomorphic} if a bijection between their ground sets
	maps one family to the other. We say that a set system $F$ is a
	\emph{configuration} of $\mathcal S$, and write $F\prec\mathcal S$, if for
	some $Y\subseteq X$, the family $F$ is isomorphic to a subfamily of
	$\mathcal S|_Y$. Otherwise we write $F\not\prec\mathcal S$.
	
	\begin{definition}
		For a fixed configuration $F$, define
		\[
		\Avoid(m,F)
		=
		\{\mathcal S\subseteq\mathcal P([m]):F\not\prec\mathcal S\},
		\]
		and
		\[
		\forb(m,F)
		=
		\max_{\mathcal S\in\Avoid(m,F)}|\mathcal S|.
		\]
	\end{definition}
	
	\subsection{The Anstee--Sali conjecture}
	
	For $n\ge 1$, define the three standard set systems on $[n]$ by
	\begin{align*}
		I_n
		&=\bigl\{\{i\}:i\in[n]\bigr\},\\
		I_n^c
		&=\bigl\{[n]\setminus\{i\}:i\in[n]\bigr\},\\
		T_n
		&=\bigl\{\varnothing,\{1\},\{1,2\},\ldots,[n]\bigr\}.
	\end{align*}
	When the block size is unimportant, we write simply $I$, $I^c$, and $T$.
	
	\begin{definition}
		Let $\mathcal A_1,\ldots,\mathcal A_p$ be set systems on pairwise disjoint
		ground sets $X_1,\ldots,X_p$. Their product is
		\[
		\mathcal A_1\times\cdots\times\mathcal A_p
		=
		\{A_1\cup\cdots\cup A_p:A_i\in\mathcal A_i\text{ for all }i\in[p]\}.
		\]
	\end{definition}
	
	For a fixed configuration $F$, let $X(F)$ be the smallest integer $p$ such
	that every $p$-fold product with factors chosen from $I$, $I^c$, and $T$
	contains $F$, provided the factor blocks are sufficiently large. Equivalently,
	$X(F)-1$ is the largest number of factors for which at least one such product
	avoids $F$.
	
	A product of $X(F)-1$ avoiding factors, taken on blocks of sizes as equal as
	possible, has order $m^{X(F)-1}$. The Anstee--Sali conjecture therefore predicts
	\[
	\forb(m,F)=\Theta\!\left(m^{X(F)-1}\right).
	\]
	
	\section{The six-vertex candidate}
	
	Let
	\[
	V=\{x,y,a,b,c,d\},
	\]
	and define
	\begin{equation}\label{eq:F2}
		\cF_2
		=
		\bigl\{
		\{x,y,a,b\},
		\{x,y,b,c\},
		\{x,y,c,d\},
		\{x,y,d,a\}
		\bigr\}.
	\end{equation}
	Thus $\cF_2$ is the $4$-uniform hypergraph obtained by adjoining the fixed
	core $\{x,y\}$ to every edge of a $4$-cycle on $\{a,b,c,d\}$.
	
	\subsection{\texorpdfstring{Computing $X(\cF_2)$}{Computing X(F2)}}
	
	\begin{proposition}\label{prop:X4}
		For the family $\cF_2$ in \eqref{eq:F2},
		\[
		X(\cF_2)=4.
		\]
	\end{proposition}
	
	\begin{proof}
		We first prove $X(\cF_2)\le 4$. Consider a four-fold product
		\[
		\mathcal A_1\times\mathcal A_2\times\mathcal A_3\times\mathcal A_4,
		\qquad
		\mathcal A_i\in\{I,I^c,T\},
		\]
		on sufficiently large, pairwise disjoint blocks.
		
		Suppose first that at least one factor is of type $I^c$ or $T$; after
		relabeling, assume it is $\mathcal A_1$. Choose distinct points
		$x,y,a$ in the first block so that the restriction of $\mathcal A_1$ contains
		both $\{x,y\}$ and $\{x,y,a\}$. This is immediate for $I^c$ by omitting $a$
		or a point outside $\{x,y,a\}$, and for $T$ by choosing $x,y,a$ in this order
		along the defining chain. Next choose one point $b,c,d$ from the second,
		third, and fourth blocks, respectively. For each of the standard factors
		$I$, $I^c$, and $T$, the restriction to a single chosen point contains both
		$\varnothing$ and that point. Hence the restriction of the four-fold product
		to $\{x,y,a,b,c,d\}$ contains
		\[
		\{x,y\}\cup B
		\qquad\text{for every }B\subseteq\{a,b,c,d\},
		\]
		and therefore contains $\cF_2$.
		
		It remains to consider the case in which all four factors are of type $I$.
		Choose $x$ from the first block, $y$ from the second block, distinct points
		$a,c$ from the third block, and distinct points $b,d$ from the fourth block.
		The restricted product then contains
		\[
		\{x,y,a,b\},\quad
		\{x,y,a,d\},\quad
		\{x,y,c,b\},\quad
		\{x,y,c,d\},
		\]
		which is exactly a copy of $\cF_2$. Thus every four-fold standard product
		contains $\cF_2$, so $X(\cF_2)\le 4$.
		
		For the reverse inequality, the three-fold product $I\times I\times I$
		avoids $\cF_2$. Every member of this product has size $3$, and restriction
		cannot increase set size, whereas every member of $\cF_2$ has size $4$.
		Therefore $X(\cF_2)>3$, completing the proof.
	\end{proof}
	
	Consequently, the Anstee--Sali conjecture predicts
	\begin{equation}\label{eq:predictionF2}
		\forb(m,\cF_2)=\Theta(m^3).
	\end{equation}
	
	\section{Mubayi's explicit construction}
	
	As a $4$-uniform hypergraph, $\cF_2$ is the complete $4$-partite
	$4$-graph
	\[
	K^{(4)}(1,1,2,2),
	\]
	with parts $\{x\}$, $\{y\}$, $\{a,c\}$, and $\{b,d\}$. The case needed
	here is the specialization $r=4$ and $t=1$ of the algebraic construction
	in the proof of Mubayi's Theorem~3.1~\cite{MUBAYI_2002}. We record the
	construction explicitly.
	
	Let $q$ be a prime power, let $\mathbb F_q$ be the field with $q$
	elements, and write $\mathbb F_q^\ast=\mathbb F_q\setminus\{0\}$. Define
	a $4$-uniform hypergraph $G_q$ with vertex set
	\[
	V(G_q)=\mathbb F_q^\ast\times\mathbb F_q^\ast.
	\]
	Thus
	\[
	n:=|V(G_q)|=(q-1)^2.
	\]
	Four distinct vertices
	\[
	(a_1,b_1),\ldots,(a_4,b_4)\in V(G_q)
	\]
	form an edge of $G_q$ precisely when
	\begin{equation}\label{eq:Mubayi-edge}
		\prod_{i=1}^4 a_i+\prod_{i=1}^4 b_i=1.
	\end{equation}
	This is exactly Mubayi's construction for $t=1$: the multiplicative
	subgroup of order $t$ is then the trivial subgroup $\{1\}$, so the
	equivalence relation appearing in the general construction is trivial.
	
	We first count the edges. Fix three distinct vertices
	$(a_i,b_i)$, $1\le i\le 3$, and put
	\[
	A=a_1a_2a_3,\qquad B=b_1b_2b_3.
	\]
	For a fourth vertex $(a,b)$, condition \eqref{eq:Mubayi-edge} becomes
	\[
	Aa+Bb=1.
	\]
	For each $a\in\mathbb F_q^\ast$ except $a=A^{-1}$, there is a unique
	nonzero choice
	\[
	b=B^{-1}(1-Aa).
	\]
	After excluding the at most three choices for which $(a,b)$ coincides
	with one of the previously chosen vertices, there are at least $q-4$
	possible fourth vertices. Since every edge is counted four times,
	according to the omitted vertex, we obtain
	\begin{equation}\label{eq:Mubayi-count}
		|E(G_q)|
		\ge
		\frac{1}{4}\binom{(q-1)^2}{3}(q-4)
		=
		\frac{1}{24}q^7-O(q^6)
		=
		\left(\frac{1}{24}+o(1)\right)n^{7/2}.
	\end{equation}
	
	We next verify directly that $G_q$ contains no copy of $\cF_2$. Suppose
	to the contrary that such a copy exists. Write its two singleton parts
	as
	\[
	(a_1,b_1),\qquad (a_2,b_2),
	\]
	and write the two parts of size two as
	\[
	(u_1,v_1),(u_2,v_2)
	\qquad\text{and}\qquad
	(x_1,y_1),(x_2,y_2).
	\]
	For $j=1,2$, set
	\[
	p_j=a_1a_2u_j,\qquad q_j=b_1b_2v_j.
	\]
	Because the two vertices $(u_1,v_1)$ and $(u_2,v_2)$ are distinct, the
	coefficient pairs $(p_1,q_1)$ and $(p_2,q_2)$ are distinct. The four
	edges of the supposed copy imply, for each $k=1,2$,
	\begin{equation}\label{eq:Mubayi-system}
		p_1x_k+q_1y_k=1,
		\qquad
		p_2x_k+q_2y_k=1.
	\end{equation}
	The two equations in \eqref{eq:Mubayi-system} have at most one solution
	$(x,y)\in(\mathbb F_q^\ast)^2$. Indeed, if their coefficient rows were
	nontrivially proportional, the two right-hand sides could not both be
	equal to $1$; otherwise the coefficient matrix is nonsingular and the
	solution is unique. Thus
	\[
	(x_1,y_1)=(x_2,y_2),
	\]
	contradicting the assumption that these are two distinct vertices.
	Therefore $G_q$ is $\cF_2$-free.
	
	Consequently, for $n=(q-1)^2$,
	\[
	\ex_4(n,\cF_2)
	\ge
	\left(\frac{1}{24}+o(1)\right)n^{7/2}.
	\]
	More generally, as in Mubayi's proof, one may choose a prime
	$q=(1-o(1))\sqrt m$ and add isolated vertices to obtain an
	$\cF_2$-free $4$-graph on $m$ vertices with $\Omega(m^{7/2})$ edges.
	This is the lower-bound half of Mubayi's asymptotic formula
	\[
	\ex_4\!\left(m,K^{(4)}(1,1,2,2)\right)
	=
	\Theta(m^{7/2}).
	\]
	
	Finally, viewing $E(G_q)$ as a set system gives an
	$\cF_2$-avoiding family. Because every member has size $4$, a
	configuration copy of the $4$-uniform family $\cF_2$ would already be an
	ordinary hypergraph copy of $\cF_2$. Hence
	\[
	\forb(m,\cF_2)
	\ge
	\ex_4(m,\cF_2)
	=
	\Omega(m^{7/2}).
	\]
	Since $7/2>3$, this contradicts the prediction in
	\eqref{eq:predictionF2}.
	
	\section{AI-assistant disclosure}
	
	The author used GPT-5.6 Sol to help identify the example and to assist with
	the exposition. The author independently reviewed and revised the outputs,
	verified the mathematical claims, and takes full responsibility for the final
	manuscript.
	
	\begin{thebibliography}{9}
		
		\bibitem{AnsteeSali2013}
		R.~P. Anstee and A.~Sali.
		A survey of forbidden configuration results.
		\emph{Electronic Journal of Combinatorics}, 20:Dynamic Survey 20, 2013.
		\newblock doi:\href{https://doi.org/10.37236/2379}{10.37236/2379}.
		
		\bibitem{MUBAYI_2002}
		D.~Mubayi.
		Some exact results and new asymptotics for hypergraph Tur\'an numbers.
		\emph{Combinatorics, Probability and Computing}, 11(3):299--309, 2002.
		\newblock doi:\href{https://doi.org/10.1017/S0963548301005028}{10.1017/S0963548301005028}.
		
		\bibitem{Sauer1972}
		N.~Sauer.
		On the density of families of sets.
		\emph{Journal of Combinatorial Theory, Series A}, 13(1):145--147, 1972.
		\newblock doi:10.1016/0097-3165(72)90019-2.
		
		\bibitem{Shelah1972}
		S.~Shelah.
		A combinatorial problem; stability and order for models and theories in
		infinitary languages.
		\emph{Pacific Journal of Mathematics}, 41(1):247--261, 1972.
		
		\bibitem{Turan1941}
		P.~Tur\'an.
		On an extremal problem in graph theory.
		\emph{Matematikai \`es Fizikai Lapok}, 48:436--452, 1941.
		
		\bibitem{VapnikChervonenkis1971}
		V.~N. Vapnik and A.~Ya. Chervonenkis.
		On the uniform convergence of relative frequencies of events to their
		probabilities.
		\emph{Theory of Probability and Its Applications}, 16(2):264--280, 1971.
		
	\end{thebibliography}
	
\end{document}
